\documentclass[pdflatex,sn-mathphys-num]{sn-jnl}

\usepackage{graphicx}
\usepackage{multirow}
\usepackage{amsmath,amssymb,amsfonts}
\usepackage{amsthm}
\usepackage{mathrsfs}
\usepackage[title]{appendix}
\usepackage{xcolor}
\usepackage{textcomp}
\usepackage{manyfoot}
\usepackage{booktabs}
\usepackage{algorithm}
\usepackage{algorithmicx}
\usepackage{algpseudocode}
\usepackage{listings}

\raggedbottom

\begin{document}

\title[Master's Thesis Proposal]{Beyond Pairwise GUI Quality Judgments: Multi-Dimensional Evaluation with Modern Multimodal LLMs}

\author{\fnm{Xinyang} \sur{Xie}}

\affil{\orgdiv{Department of Informatics}, \orgname{Technische Universitaet Clausthal},
\orgaddress{\city{Clausthal-Zellerfeld}, \country{Germany}}}

\abstract{
Assessing the quality of graphical user interfaces (GUIs) is becoming increasingly
important as large language models are used to generate interface designs and frontend
code. Existing screenshot-based approaches can support GUI quality judgments, but they
remain limited by binary decisions, older models, and simple prompting strategies.
This proposal extends that line of work with current multimodal large language models.
The study combines pairwise comparison with rubric-based assessment across layout
quality, visual hierarchy, information organization, and perceived usability.
Methodologically, it uses a screenshot-based evaluation set, a human annotation
protocol with agreement analysis, and a comparison of judging strategies such as
zero-shot prompting, rubric-guided scoring, repeated judging, and order-swapped
pairwise comparisons, where the order of candidates is reversed to detect positional
bias. The resulting protocol will be packaged into a benchmark and lightweight
leaderboard that can be extended to newly added models. In addition to the static
screenshot-based track, the proposal includes a small dynamic evaluation module in
which computer-use agents attempt short tasks on selected GUIs. The goal is to
determine which evaluation setup aligns best with human judgments and how static
quality judgments relate to dynamic interaction performance.
}

\maketitle

%%------------------------------------------------------------
\section{Introduction and Motivation}\label{sec:intro}
%%------------------------------------------------------------

The quality of a graphical user interface (GUI) shapes how efficiently users complete
tasks and how they perceive a software system. As generative models are increasingly
used to produce interface mockups and frontend code, the question of how UI quality can
be evaluated becomes more pressing.

UIClip~\cite{uiclip} provided an important first step by showing that vision-language
models can perform pairwise comparisons of UI screenshots and identify which of two
interfaces appears better. At the same time, the original setup leaves several open
questions: it relies on older models, reduces quality assessment to binary comparisons,
and does not examine how newer LLM-as-a-judge strategies affect reliability.

Recent work on improving UI generation from designer feedback~\cite{designer_feedback}
suggests that high-quality, workflow-aligned feedback can be more informative than
large amounts of coarse ranking data. In parallel, the broader LLM-as-a-judge
literature~\cite{llm_judge} has introduced prompting, consistency, and aggregation
strategies that may also be useful for GUI evaluation. Taken together, these
developments motivate a more methodologically explicit re-evaluation of GUI quality
judgment with contemporary multimodal models.

At the same time, screenshot-based assessment captures only part of GUI quality. It
can reveal visual structure and perceived usability, but it cannot show whether an
interface supports successful interaction once users or agents start navigating it. For
this reason, the proposal treats dynamic task completion as a complementary evaluation
dimension rather than relying on static screenshots alone.

This thesis therefore asks whether modern multimodal large language models can
approximate human judgments of GUI quality beyond pairwise comparison, and how such
judgments should be elicited and evaluated in a reliable way.

%%------------------------------------------------------------
\section{Planned Approach}\label{sec:approach}
%%------------------------------------------------------------

The planned work centers on an evaluation pipeline with six main components: dataset
construction, human annotation, rubric design, the comparison of LLM-as-a-judge
strategies, benchmark and leaderboard design, and a small dynamic evaluation
extension.

\subsection{Datasets and Evaluation Material}

The study will build primarily on publicly available screenshot-based resources from
UIClip. To avoid evaluating only older interfaces, this material will be complemented,
if needed, by a smaller curated set of contemporary web or mobile screenshots. The
resulting evaluation material will consist of two subsets: a pairwise comparison subset
for overall preference judgments and a single-image subset for structured
multi-dimensional assessment.

\subsection{Human Annotation Protocol}

Human judgments will be collected for a selected evaluation subset. For pairwise tasks,
annotators will compare two GUI screenshots and indicate which interface is better
overall, with an optional tie label if required. For single-image tasks, annotators
will score each GUI with a small rubric and provide brief justifications. To improve
the reliability of the reference labels, each item will be annotated by multiple
raters, and inter-rater reliability will be measured, for example with
Krippendorff's alpha.

\subsection{Multi-dimensional GUI Rubric}

To move beyond binary comparison, I will define a small rubric grounded in common
visual design principles and usability heuristics. The initial dimensions are:

\begin{itemize}
    \item \textbf{Layout quality}: spatial organization, alignment, spacing, and
    overall balance.
    \item \textbf{Visual hierarchy}: whether contrast, size, and color guide attention
    appropriately.
    \item \textbf{Information organization}: clarity of grouping and logical ordering
    of content.
    \item \textbf{Perceived usability}: how easily key actions and interface purpose
    can be inferred from a static screenshot.
\end{itemize}

Each dimension will be scored on a fixed ordinal scale with short written anchors so
that the same rubric can be applied by both human annotators and model judges. The
perceived usability dimension is intentionally limited to what can reasonably be
inferred from a static screen, rather than to full interaction quality.

\subsection{Prompting and Judging Strategies}

Several judging strategies will be compared. A zero-shot baseline will ask the model
directly for a pairwise preference or a rubric-based score. A few-shot variant will add
one or two worked examples to test whether light calibration improves judgment quality.
A rubric-guided condition will provide explicit scoring instructions and require a
structured output format. To test robustness, the same prompt will be repeated to
measure self-consistency. For pairwise tasks, order-swapped pairwise comparisons will
be used, where the order of candidates is reversed to detect positional bias. In
addition, the study will compare single-model judgments with simple aggregation
strategies such as majority voting for pairwise tasks and averaged scores for
rubric-based tasks.

\subsection{Benchmark and Leaderboard Design}

Beyond comparing individual models, the project will package the evaluation setup into
a benchmark structure that combines pairwise tasks, rubric-based scoring, and shared
reference annotations. Results will be stored in a standardized format so that newly
added multimodal models can be evaluated under the same protocol. As a practical
artifact, the benchmark will be accompanied by a lightweight leaderboard or web-based
dashboard for comparing models across tasks and rubric dimensions.

\subsection{Dynamic Evaluation with Computer-Use Agents}

To complement the static benchmark, the project will include a small dynamic evaluation
extension. For a selected subset of GUIs, short interaction tasks will be defined and
executed by computer-use agents in a browser-based environment. This extension will be
kept deliberately small and will focus on whether interfaces that receive strong static
quality judgments also support reliable task completion in practice.

%%------------------------------------------------------------
\section{Evaluation}\label{sec:evaluation}
%%------------------------------------------------------------

Several current multimodal models will be tested as judges, including representative
models from major closed-source and open multimodal model families, depending on
availability at evaluation time. The pairwise setting will be evaluated through
agreement with human majority labels, while the rubric-based setting will be evaluated
through rank correlation and score differences between model outputs and human ratings.
The dynamic extension will be evaluated through task completion rate, action
efficiency, and qualitative failure analysis on the selected interactive tasks.

Beyond task accuracy, the evaluation will explicitly examine reliability. This includes
inter-rater agreement in the human annotations, consistency across repeated model
judgments, stability under order-swapped pairwise comparisons, and agreement across
different models. Computational cost will also be tracked in order to assess practical
trade-offs between evaluation quality and API usage. A further goal is to determine
whether the resulting benchmark and leaderboard provide a stable and extensible basis
for comparing both current and newly introduced multimodal models.

Finally, a qualitative error analysis will examine typical failure cases, such as
cluttered layouts, minimal interfaces, text-heavy screens, and visually appealing but
less usable designs. Together, these analyses should clarify how well modern
multimodal large language models support UIClip-style evaluation, which rubric
dimensions are most reliable, and which judging strategies align best with human
assessment.

\backmatter

{\footnotesize\setlength{\bibsep}{0pt}\bibliography{references}}

\end{document}
